\documentclass[../main.tex]{subfiles}
\begin{document}

The pruning stage of Section~\ref{sec:stage1} ranks every node by the combined
importance of Eq.~\eqref{eq:importance},
$\mathrm{Imp}_v = \widehat{\mathrm{Inf}}_v^{\,\alpha}\,\widehat{\mathrm{Rel}}_v^{\,1-\alpha}$,
where the blend coefficient $\alpha \in [0,1]$ sets how much the output-side
\emph{influence} and the input-side \emph{relevance} each contribute. The two
endpoints recover familiar criteria: $\alpha = 1$ keeps nodes by influence alone,
which is the pruning rule of the original circuit-tracing work of Ameisen et
al.~\cite{ameisen2025circuit}, while $\alpha = 0$ keeps them by relevance alone.
This appendix isolates the effect of $\alpha$, so that the gain from adding
relevance to the influence-only baseline can be read off directly.

\section{Experimental setup}
The ablation runs on the $100$ pruned attribution graphs of the analogies dataset
(Section~\ref{sec:datasets}). For each graph we sweep the blend coefficient
$\alpha \in \{0, 0.25, 0.5, 0.75, 1\}$ against the node threshold
$\eta_V \in \{0.01, 0.02, 0.05, 0.1, 0.2, 0.4, 0.6, 0.8, 1.0\}$, holding the
remaining knobs fixed at the defaults used elsewhere in the thesis: the geometric
combiner, rank normalization of the two scores, an edge threshold of $0.95$, and
target-logit weighting. Because $\eta_V$ is a retained-mass fraction rather than a
node count, the same $\eta_V$ yields a different subgraph size for each $\alpha$;
we therefore report every metric against the resulting \emph{average node count}
on a logarithmic axis, which makes the curves comparable at equal graph size. The
relevance score reads the input tokens through SHAP token weights, and these
weights depend on a normalization choice; we run the whole sweep under both
normalizations used in the thesis, entmax-$1.25$ and softmax, to confirm the
conclusion does not hinge on that choice.

\section{Faithfulness metrics}
We grade each pruned subgraph with three metrics, all oriented so that higher is
better. The \emph{replacement score} and the \emph{graph completeness score} are
the cross-layer-transcoder graph scores of Ameisen et
al.~\cite{ameisen2025circuit}: every pruned feature is folded into the
reconstruction-error node at its own layer and position, and the scores then
measure the fraction of logit-directed influence that the retained features carry
rather than leak into those error nodes. Both are \emph{output-side} measures, as
they weight nodes by their influence on the target logit. The third metric,
\emph{token-attribution faithfulness}, is instead an \emph{input-side} measure: for
each input token it computes the relevance that reaches the target logit through
the pruned graph as a fraction of the relevance that reaches it through the full
graph, and averages this ratio over tokens weighted by their SHAP attribution. The
split is deliberate: the first two metrics reward reconstructing the output-side
computation that influence targets, whereas the third rewards preserving the
input-side attribution that relevance targets, so together they probe both ends of
the pruning criterion.

\section{Effect of the blend coefficient}
Figure~\ref{fig:prune-alpha-sweep} reports the three metrics against average node
count for every $\alpha$, under both token-weight normalizations. The two
normalizations produce nearly indistinguishable curves, so the discussion below
refers to the entmax panels of Figure~\ref{fig:prune-alpha-entmax} and applies
equally to the softmax panels of Figure~\ref{fig:prune-alpha-softmax}.

\begin{figure}[htbp]
\centering
\begin{subfigure}{\linewidth}
\centering
\includegraphics[width=\linewidth]{Figure/prune_alpha_sweep_entmax.png}
\caption{entmax-$1.25$ SHAP token weighting.}
\label{fig:prune-alpha-entmax}
\end{subfigure}

\vspace{0.8em}

\begin{subfigure}{\linewidth}
\centering
\includegraphics[width=\linewidth]{Figure/prune_alpha_sweep_softmax.png}
\caption{softmax SHAP token weighting.}
\label{fig:prune-alpha-softmax}
\end{subfigure}
\caption[Influence--relevance blend ablation for pruning]{Three faithfulness
metrics --- graph completeness score, replacement score, and token-attribution
faithfulness --- against average node count (log axis) for blend coefficients
$\alpha \in \{0, 0.25, 0.5, 0.75, 1\}$, averaged over the $100$ analogy graphs.
$\alpha = 1$ is influence-only pruning (the original circuit-tracing
criterion~\cite{ameisen2025circuit}) and $\alpha = 0$ is relevance-only pruning.
Each curve sweeps the node threshold $\eta_V$ from $0.01$ (left, most aggressive)
to $1.0$ (right, full graph).}
\label{fig:prune-alpha-sweep}
\end{figure}

The two output-side metrics rank the blends as expected. On both the completeness
and the replacement score, a larger $\alpha$ gives a higher score at any fixed node
budget, and the influence-only setting $\alpha = 1$ is the strongest once the graph
is allowed to keep a few hundred nodes; at $\eta_V = 0.2$ (about $570$ nodes) it
reaches a completeness of $0.82$ and a replacement of $0.63$, against $0.65$ and
$0.45$ for relevance-only pruning. This ordering is the one the metrics are built
to produce, since both reward reconstructing the logit-side influence that the
influence score targets directly.

The token-attribution faithfulness metric reverses that ordering. Here the
relevance-weighted blends dominate the influence-only baseline across the entire
mid-range: at the same $\eta_V = 0.2$ operating point, relevance-only pruning
attains a token-attribution faithfulness of $0.41$ while influence-only reaches
only $0.14$, and every intermediate $\alpha$ falls between the two. The reversal is
expected for the same structural reason as before: this metric measures how well
the pruned graph preserves the attribution flowing from the input tokens, which is
exactly the quantity the relevance score is built from, so dropping relevance from
the criterion costs the most here.

The regime that matters for compact, human-readable circuits is the aggressive left
end of the curves, and there the influence-only advantage on the output-side
metrics disappears. Table~\ref{tab:prune-alpha-aggressive} reads off the most
aggressive threshold $\eta_V = 0.01$, where each graph keeps about $25$ nodes. At
this operating point the relevance-aware blends are not merely competitive with
influence-only pruning on the output-side metrics --- they match or slightly exceed
it: the $\alpha = 0.75$ blend reaches a completeness of $0.567$ and a replacement of
$0.413$ against the influence-only $0.549$ and $0.408$, while every intermediate
blend simultaneously leads on token-attribution faithfulness ($0.086$--$0.089$
versus $0.067$). Only pure relevance ($\alpha = 0$) trails on the output-side
scores. The
influence-only baseline pulls ahead on completeness and replacement only as the node
budget grows past roughly a hundred nodes, a region already too large for the
interpretable summaries the pipeline targets.

\begin{table}[h]
\centering
\small
\caption{Pruning faithfulness at the most aggressive threshold
($\eta_V = 0.01$, about $25$ kept nodes on average), entmax token weighting, mean
over the $100$ analogy graphs. At this operating point the influence--relevance
blends ($0 < \alpha < 1$) match or exceed the influence-only baseline
($\alpha = 1$) on all three metrics; bold marks the best value per column. The
softmax sweep gives the same ordering.}
\label{tab:prune-alpha-aggressive}
\begin{tabular}{l c c c}
\toprule
$\alpha$ & completeness $\uparrow$ & replacement $\uparrow$ & token-attr.\ faithfulness $\uparrow$ \\
\midrule
$0.00$ (relevance only) & 0.470 & 0.381 & 0.061 \\
$0.25$                  & 0.549 & 0.405 & \textbf{0.089} \\
$0.50$                  & 0.559 & 0.409 & 0.087 \\
$0.75$                  & \textbf{0.567} & \textbf{0.413} & 0.086 \\
$1.00$ (influence only) & 0.549 & 0.408 & 0.067 \\
\bottomrule
\end{tabular}
\end{table}

Taken together, the sweep shows that adding relevance to the pruning criterion costs
essentially nothing in output-side faithfulness at the aggressive operating point
the pipeline uses, while strictly improving the input-side token-attribution
faithfulness. The output-side metrics, however, give no credit for keeping
semantically meaningful, input-relevant nodes, which is the property relevance is
meant to protect. That complementary benefit is assessed separately by the
relevance-oriented evaluation of Section~\ref{sec:llm-results}, and it is the reason
the pipeline retains a non-zero relevance weight rather than defaulting to the
influence-only criterion.

\end{document}
